\documentclass[11pt,a4paper]{article}
\usepackage[margin=1in]{geometry}
\usepackage{times}
\usepackage{latexsym}
\usepackage{amsmath}
\usepackage{amssymb}
\usepackage{booktabs}
\usepackage{multirow}
\usepackage{graphicx}
\usepackage{xcolor}
\usepackage{url}
\usepackage{hyperref}
\usepackage{microtype}
\usepackage{subcaption}
\usepackage{cite}

\newcommand{\gap}{\textsc{EAGap}}
\newcommand{\ratio}{\ensuremath{r_{\text{sel}}}}

\title{The Embedding Alignment Gap: Internal Diversity Metrics Diverge\\from Content Diversity in Iterative Self-Training}

\author{Anonymous}

\begin{document}
\maketitle

\begin{abstract}
Iterative self-training---where a model is repeatedly fine-tuned on its own generations---is known to cause model collapse.
A natural mitigation is diversity-based data selection using the model's internal embeddings.
We show that this strategy backfires: a systematic \emph{embedding alignment gap} emerges in which the geometric diversity of internal representations diverges from content diversity (measured by rare $n$-gram survival) across generations.
Concretely, the internal \emph{selection ratio}---the fraction of content-diverse samples among those deemed geometrically diverse---declines from ${\sim}1.03$ to ${\sim}0.78$ over 8 generations of GPT-2 self-training ($p{=}0.003$ across 3 seeds), while frozen external embeddings (Contriever) maintain a ratio of ${\sim}0.99$.
Internal diversity selection increases perplexity by 13.0\% relative to uncurated self-training ($p{=}0.02$), \emph{failing to prevent and potentially worsening} collapse.
Two measurement dissociation phenomena underlie these results: (i) perplexity and rare $n$-gram frequency decouple---external selection improves perplexity by 3.8\% without improving lexical diversity; (ii) RankMe, a standard representation quality metric, \emph{increases} during collapse despite continuous rare $n$-gram decline.
These findings hold across two model scales (GPT-2 125M, 355M) and two datasets (WikiText-103, OpenWebText).
We argue that embedding-based diversity selection is unreliable in iterative settings and that the field needs content-aware alternatives.
\end{abstract}

\section{Introduction}

Iterative self-training, in which a language model is fine-tuned on data generated by a previous version of itself, has emerged as a common paradigm in synthetic data augmentation and knowledge distillation \cite{shumailov2024model,alemohammad2024selfconsuming}.
A well-documented failure mode is \emph{model collapse}: successive generations exhibit declining output diversity, loss of rare tokens, and degraded perplexity \cite{dohmatob2024tale}.

A natural countermeasure is diversity-based data selection.
Methods such as SemDeDup \cite{abbas2023semdedup} and D4 \cite{tirumala2023d4} use embedding-space geometry to select maximally diverse subsets.
An intuitive extension is to apply these techniques within the iterative loop, using the model's own representations to curate each generation's training data.
We investigate whether this approach prevents collapse.

Our central finding is negative: \textbf{internal-embedding diversity selection accelerates collapse rather than preventing it.}
We trace this to what we call the \emph{embedding alignment gap} (\gap{})---a progressive divergence between the model's geometric notion of diversity and the actual content diversity of its outputs.
As training proceeds, the embedding space contracts and reorients so that samples the model deems ``diverse'' are increasingly vocabulary-poor.
The selection process then amplifies homogeneity rather than counteracting it.

This paper makes three contributions:
\begin{enumerate}
\item \textbf{The embedding alignment gap.}
We define and measure a \emph{selection ratio} $\ratio$ that tracks how well embedding-based diversity selection preserves content-diverse samples.
Across 3 seeds, $\ratio$ exhibits an overall declining trend from ${\sim}1.03$ to $0.78 \pm 0.02$ over 8 generations ($p{=}0.003$), with transient recoveries at different generational positions across seeds.
Frozen external embeddings (Contriever) maintain $\ratio \geq 0.97$ throughout.

\item \textbf{Measurement dissociation.}
We identify two cases where standard metrics fail to capture diversity collapse:
(a) perplexity and rare $n$-gram frequency decouple---external selection improves perplexity without improving lexical diversity;
(b) RankMe \cite{garrido2023rankme}, a representation quality metric, \emph{increases} during the first generations of collapse before slowly declining, while rare $n$-grams decline monotonically.

\item \textbf{Cross-condition robustness.}
The alignment gap is consistent across model scales (GPT-2 125M and 355M), datasets (WikiText-103 and OpenWebText), and random seeds (3 seeds for 125M, 2 for 355M).
\end{enumerate}

\section{Background and Related Work}

\paragraph{Model collapse in iterative training.}
\cite{shumailov2024model} showed that iteratively training on model-generated data leads to progressive loss of distributional tails.
\cite{alemohammad2024selfconsuming} formalized this as ``self-consuming generative models'' and established convergence-to-delta results.
\cite{dohmatob2024tale} provided finite-sample analysis showing tail trimming accelerates with generation count.
Our work extends this line by showing that diversity-based curation, a seemingly natural countermeasure, can worsen the problem.

\paragraph{Embedding-based data selection.}
SemDeDup \cite{abbas2023semdedup} removes near-duplicates in embedding space; D4 \cite{tirumala2023d4} selects diverse subsets via facility location; DSIR \cite{xie2023dsir} reweights by importance.
These methods operate on \emph{static} datasets with \emph{frozen} embeddings.
We show that the iterative setting introduces a failure mode absent from static curation: the embeddings used for selection drift as the model trains, creating a feedback loop.

\paragraph{Neural collapse and representation geometry.}
\cite{papyan2020prevalence} showed that deep network representations converge to a simplex ETF structure during training, reducing intra-class variability.
In our setting, iterative self-training amplifies this effect: the embedding space contracts over generations, collapsing the geometric signal that diversity selection relies on.

\paragraph{Representation quality metrics.}
RankMe \cite{garrido2023rankme} uses the entropy of singular values to estimate effective dimensionality.
We find that RankMe can \emph{increase} during early collapse, revealing a blind spot: it tracks distributional spread of representations but not alignment with content-level diversity.

\section{Experimental Setup}

\subsection{Iterative Self-Training Protocol}

We train GPT-2 125M and GPT-2 355M (medium) on WikiText-103 \cite{merity2016pointer} and an OpenWebText subset (50K documents).
Starting from pretrained weights (Gen 0), each generation:
(1) generates a synthetic corpus from the current model,
(2) optionally applies a data selection method, and
(3) fine-tunes a fresh copy of the pretrained model on the selected data.
We iterate for 8 generations (Gen 0--7).

\subsection{Selection Methods}

We compare three conditions:
\begin{itemize}
\item \textbf{Uncurated}: All synthetic data is used without selection (baseline).
\item \textbf{RGDS} (Representation-Guided Diversity Selection): A submodular facility-location objective in the model's \emph{own} embedding space selects the top-$k$ most diverse samples. Embeddings are extracted from the penultimate layer of the current-generation model.
\item \textbf{External-only}: Identical selection procedure, but using frozen Contriever \cite{izacard2022contriever} embeddings instead of the model's internal representations.
\end{itemize}

\subsection{Metrics}

\paragraph{Selection ratio ($\ratio$).}
For each generation, we compute the fraction of content-diverse samples (top quartile by rare $n$-gram frequency) among those selected by the diversity method, normalized by the expected fraction under random selection.
$\ratio > 1$ indicates the selection method preferentially retains content-diverse samples; $\ratio < 1$ indicates it preferentially discards them.

\paragraph{Content diversity.}
We measure rare $n$-gram frequency: the fraction of unique 4-grams in generated text that appear $\leq 5$ times in the original training corpus.

\paragraph{Model quality.}
Perplexity (PPL) on a held-out validation set measures language modeling quality.

\paragraph{Representation quality.}
RankMe \cite{garrido2023rankme} measures effective dimensionality of the embedding space.
$k$-NN probe accuracy on a held-out classification task measures representation discriminability.
Spearman rank correlation ($\rho$) between embedding-space diversity rank and content diversity rank provides a direct alignment measure.

\section{Results}

\subsection{Main Experiment: Three-Method Comparison}
\label{sec:main}

Table~\ref{tab:main} presents the full generational trajectory for GPT-2 125M on WikiText-103 (seed=42).

\begin{table*}[t]
\centering
\small
\begin{tabular}{l ccc ccc ccc cc}
\toprule
& \multicolumn{3}{c}{Perplexity ($\downarrow$)} & \multicolumn{3}{c}{Rare $n$-gram ($\uparrow$)} & \multicolumn{3}{c}{Selection Ratio} & \multicolumn{2}{c}{Rank Corr.} \\
\cmidrule(lr){2-4} \cmidrule(lr){5-7} \cmidrule(lr){8-10} \cmidrule(lr){11-12}
Gen & Unc. & RGDS & Ext. & Unc. & RGDS & Ext. & RGDS & Ext. & $\Delta$ & RGDS & Ext. \\
\midrule
0 & 37.0 & 37.2 & 36.9 & .029 & .029 & .029 & 1.03 & 1.05 & $-$0.02 & $-$.024 & $-$.024 \\
1 & 44.8 & 41.8 & 43.2 & .019 & .018 & .019 & 0.98 & 1.06 & $-$0.08 & $-$.042 & --- \\
2 & 51.5 & 46.5 & 48.8 & .014 & .012 & .014 & 0.93 & 1.04 & $-$0.11 & $-$.056 & --- \\
3 & 56.8 & 50.1 & 52.5 & .012 & .009 & .011 & 0.87 & 1.03 & $-$0.16 & $-$.069 & --- \\
4$^\dagger$ & 55.5 & 51.8 & 51.4 & .010 & .008 & .010 & 0.91 & 1.02 & $-$0.11 & $-$.063 & --- \\
5 & 51.6 & 53.2 & 48.9 & .010 & .007 & .009 & 0.84 & 1.01 & $-$0.17 & $-$.072 & --- \\
6 & 48.5 & 53.8 & 46.5 & .009 & .007 & .008 & 0.80 & 1.00 & $-$0.20 & $-$.081 & --- \\
7 & 48.0 & 54.2 & 46.2 & .009 & .006 & .008 & 0.76 & 0.99 & $-$0.23 & $-$.090 & $-$.053 \\
\midrule
$\Delta$(0$\to$7) & +11.0 & +17.0 & +9.3 & $-$.020 & $-$.023 & $-$.021 & $-$0.27 & $-$0.06 & & $-$.066 & $-$.029 \\
\bottomrule
\end{tabular}
\caption{Generational trajectory for GPT-2 125M on WikiText-103 (seed=42). RGDS uses the model's own embeddings for diversity selection; External uses frozen Contriever embeddings. $^\dagger$Gen 4 shows a transient rebound in RGDS selection ratio. Unc.\ = Uncurated baseline.}
\label{tab:main}
\end{table*}

Three patterns emerge.
First, the RGDS selection ratio exhibits an \textbf{overall declining trend} from 1.03 to 0.76 ($\Delta{=}{-}0.27$), with a transient recovery at Gen~4 ($0.87 \to 0.91$).
This rebound coincides with a PPL plateau in the uncurated condition (Gen 3--4), suggesting a brief decompression of the embedding space when the learning signal temporarily stabilizes.
Second, the external selection ratio remains close to 1.0 throughout (1.05$\to$0.99), confirming that frozen embeddings are largely immune to the alignment gap.
Third, by Gen~7, RGDS \emph{worsens} perplexity by 12.9\% relative to uncurated self-training ($54.2$ vs $48.0$), while external selection \emph{improves} it by 3.8\% ($46.2$ vs $48.0$).

Table~\ref{tab:derived} summarizes the derived comparison metrics at Gen~7.

\begin{table}[t]
\centering
\small
\begin{tabular}{lccc}
\toprule
Metric & RGDS & Ext. & Unc. \\
\midrule
PPL (Gen 7) & 54.2 & 46.2 & 48.0 \\
PPL excess vs Unc. & +12.9\% & $-$3.8\% & --- \\
Rare $n$-gram (Gen 7) & .006 & .008 & .009 \\
Rare deficit vs Unc. & $-$33.3\% & $-$11.1\% & --- \\
Selection ratio (Gen 7) & 0.76 & 0.99 & --- \\
Rank corr. (Gen 7) & $-$.090 & $-$.053 & $-$.085 \\
\bottomrule
\end{tabular}
\caption{Derived metrics at Gen 7 (GPT-2 125M, WikiText-103, seed=42). RGDS exhibits worse PPL and lower rare $n$-gram frequency than even uncurated self-training.}
\label{tab:derived}
\end{table}

\subsection{Cross-Seed Robustness}
\label{sec:seeds}

We replicate the full experiment with seeds 123 and 456 (Table~\ref{tab:crossseed}).
Across 3 seeds, the RGDS selection ratio at Gen~7 averages $0.78 \pm 0.02$ while external averages $0.99 \pm 0.01$ (paired $t$-test $p{=}0.003$).
PPL excess of RGDS over uncurated averages $13.0 \pm 0.5\%$ ($p{=}0.02$).
Rare $n$-gram deficit averages $-30\%$ ($p{=}0.04$).

Each seed exhibits 1--2 transient rebounds in the RGDS selection ratio, but at different generational positions (Gen~4, 5, and 6 for seeds 42, 123, and 456 respectively).
This seed-dependent variability in rebound timing, combined with the consistent overall declining trend, supports describing the trajectory as an \emph{overall declining trend with transient recovery} rather than a monotonic degradation.

\begin{table}[t]
\centering
\small
\begin{tabular}{lcc}
\toprule
Metric (Gen 7) & Mean $\pm$ Std & $p$ \\
\midrule
RGDS $\ratio$ & $0.78 \pm 0.02$ & \multirow{2}{*}{\textbf{0.003}} \\
External $\ratio$ & $0.99 \pm 0.01$ & \\
\midrule
RGDS PPL excess (\%) & $13.0 \pm 0.5$ & \textbf{0.02} \\
\midrule
RGDS rare $n$-gram & $.0063 \pm .0006$ & \multirow{2}{*}{\textbf{0.04}} \\
Unc.\ rare $n$-gram & $.0087 \pm .0005$ & \\
\midrule
Rank corr.\ $\Delta$ (int.) & $-0.060 \pm 0.008$ & --- \\
\bottomrule
\end{tabular}
\caption{Cross-seed summary (GPT-2 125M, WikiText-103, 3 seeds). Bold $p$-values indicate significance at $\alpha{=}0.05$.}
\label{tab:crossseed}
\end{table}

\subsection{Model Scale: GPT-2 355M}
\label{sec:scale}

To address the concern that GPT-2 125M is too small to be informative, we replicate with GPT-2 medium (355M) on WikiText-103.
Table~\ref{tab:scale} shows the cross-scale comparison.

\begin{table}[t]
\centering
\small
\begin{tabular}{lcccc}
\toprule
& \multicolumn{2}{c}{GPT-2 125M} & \multicolumn{2}{c}{GPT-2 355M} \\
\cmidrule(lr){2-3} \cmidrule(lr){4-5}
Metric (Gen 7) & $s{=}42$ & $s{=}123$ & $s{=}42$ & $s{=}123$ \\
\midrule
RGDS $\ratio$ & 0.76 & 0.78 & 0.80 & 0.82 \\
Ext.\ $\ratio$ & 0.99 & 0.98 & 0.98 & 0.97 \\
RGDS PPL & 54.2 & 55.1 & 43.2 & 44.0 \\
Unc.\ PPL & 48.0 & 48.5 & 37.8 & 38.5 \\
PPL excess (\%) & 12.9 & 13.6 & 14.3 & 14.3 \\
RGDS rare & .006 & .006 & .007 & .007 \\
Unc.\ rare & .009 & .008 & .010 & .010 \\
Rare deficit (\%) & $-$33 & $-$25 & $-$30 & $-$30 \\
Rebounds & Gen 4 & Gen 5 & none & Gen 5 \\
\bottomrule
\end{tabular}
\caption{Cross-scale, cross-seed comparison. The alignment gap persists at 355M parameters, with consistent PPL excess (${\sim}13$--$14\%$) and rare $n$-gram deficit (${\sim}25$--$33\%$). 355M shows a smoother trajectory (fewer rebounds) and slightly higher terminal $\ratio$.}
\label{tab:scale}
\end{table}

The alignment gap is present at both scales.
GPT-2 355M shows a slightly higher terminal $\ratio$ ($0.81$ mean vs $0.77$ for 125M, difference ${\sim}0.04$), suggesting that larger model capacity allows a somewhat slower embedding space contraction.
However, the PPL excess is comparable ($14.3\%$ vs $13.0\%$) and rare $n$-gram deficit is similar ($-30\%$ vs $-30\%$), indicating that the downstream consequence of the gap is scale-invariant in this range.
Notably, GPT-2 355M exhibits stronger self-recovery in the uncurated condition (PPL rises to 43.2 at Gen~3 then falls to 37.8 by Gen~7), yet RGDS still accumulates a 14.3\% PPL excess, confirming that internal selection actively interferes with this recovery mechanism.

\subsection{Cross-Dataset Robustness: OpenWebText}
\label{sec:owt}

Table~\ref{tab:dataset} compares WikiText-103 and OpenWebText (GPT-2 125M, seed=42).

\begin{table}[t]
\centering
\small
\begin{tabular}{lcc}
\toprule
Metric (Gen 7) & WikiText-103 & OpenWebText \\
\midrule
RGDS $\ratio$ & 0.76 & 0.74 \\
External $\ratio$ & 0.99 & 1.01 \\
RGDS PPL excess (\%) & 12.9 & 14.3 \\
Rare deficit (\%) & $-$33 & $-$29 \\
Rank corr.\ (RGDS) & $-$.090 & $-$.095 \\
Rank corr.\ (Ext.) & $-$.053 & $-$.050 \\
\bottomrule
\end{tabular}
\caption{Cross-dataset comparison (GPT-2 125M, seed=42). Direction and magnitude of the alignment gap are consistent across WikiText-103 and OpenWebText.}
\label{tab:dataset}
\end{table}

The alignment gap manifests consistently across both datasets.
On OpenWebText, the RGDS selection ratio is slightly lower (0.74 vs 0.76) and the PPL excess slightly higher (14.3\% vs 12.9\%), which may reflect that OpenWebText's more homogeneous text distribution (web crawl vs.\ curated Wikipedia) leads to faster embedding space compression.
Importantly, the RGDS trajectory on OpenWebText shows no rebounds, unlike WikiText-103, further supporting the interpretation that rebounds are data/seed-dependent noise rather than a systematic phenomenon.

\section{Measurement Dissociation Analysis}
\label{sec:dissociation}

Beyond the alignment gap itself, our experiments reveal two measurement dissociation phenomena that have broader implications for monitoring model collapse.

\subsection{Perplexity--Lexical Diversity Decoupling}
\label{sec:ppl-rare}

External selection improves perplexity by 3.8\% relative to uncurated training ($46.2$ vs $48.0$ at Gen~7) but does \emph{not} improve rare $n$-gram survival ($0.008$ vs $0.009$, a deficit of $11.1\%$).
This decoupling is consistent across datasets: on OpenWebText, external selection achieves PPL of 36.8 (vs.\ 38.5 uncurated, $-4.4\%$) while rare $n$-gram frequency is 0.007 (vs.\ 0.007 uncurated).

This dissociation means that semantic diversity (as captured by Contriever embeddings) and lexical diversity (as measured by rare $n$-gram frequency) are distinct phenomena under iterative self-training.
A method can select data that is semantically diverse---covering different topics, styles, and registers---while still failing to preserve the lexical tail.
Rare tokens are lost not because the selected data lacks topical variety, but because the generation process itself has stopped producing them.

\subsection{RankMe Fails to Detect Diversity Collapse}
\label{sec:rankme}

RankMe, which estimates the effective dimensionality of the representation space via singular value entropy, shows a striking dissociation from content diversity.
In the RGDS condition (seed=42, GPT-2 125M), RankMe \emph{increases} from 281 at Gen~0 to a peak of 305 at Gen~3, before gradually declining to 292 at Gen~7.
Simultaneously, rare $n$-gram frequency drops monotonically from 0.029 to 0.006---an 79\% decline.

The initial RankMe increase likely reflects the early-generation expansion of the embedding manifold as the model encounters diverse synthetic data.
But this expansion is geometric, not content-aligned: the model spreads representations across the embedding space in directions uncorrelated with lexical diversity.
By the time RankMe begins declining (Gen~4+), content diversity has already been severely degraded.

Similarly, $k$-NN probe accuracy remains at 0.988 from Gen~0 through Gen~7, completely insensitive to the underlying collapse.
These findings caution against using representation quality metrics as proxies for output diversity in iterative settings.

Table~\ref{tab:rankme} quantifies the RankMe--rare dissociation.

\begin{table}[t]
\centering
\small
\begin{tabular}{lcccc}
\toprule
Gen & RankMe & $\Delta$RankMe & Rare & $\Delta$Rare \\
\midrule
0 & 281 & --- & .029 & --- \\
1 & 296 & +15 & .018 & $-$.011 \\
2 & 301 & +5 & .012 & $-$.006 \\
3 & 305 & +4 & .009 & $-$.003 \\
4 & 303 & $-$2 & .008 & $-$.001 \\
5 & 299 & $-$4 & .007 & $-$.001 \\
6 & 295 & $-$4 & .007 & $-$.001 \\
7 & 292 & $-$3 & .006 & $-$.001 \\
\bottomrule
\end{tabular}
\caption{RankMe vs.\ rare $n$-gram frequency across generations (RGDS, GPT-2 125M, seed=42). RankMe increases through Gen~3 while rare $n$-grams decline monotonically, revealing a measurement dissociation.}
\label{tab:rankme}
\end{table}

\section{Discussion}

\paragraph{Why does the gap emerge?}
We hypothesize that iterative self-training induces a form of progressive neural collapse \cite{papyan2020prevalence} in the representation space.
As the model trains on increasingly homogeneous data, intra-class variability in the embedding space decreases.
The submodular diversity objective, operating in this contracted space, selects points that are geometrically spread along the remaining principal components---but these components are no longer aligned with content diversity dimensions.
The selection increasingly favors samples that vary along collapse-mode directions (e.g., minor syntactic variations) rather than content-relevant dimensions (e.g., rare vocabulary usage).

\paragraph{Transient rebounds.}
The 1--2 generation rebounds in RGDS selection ratio (occurring at different positions across seeds) are consistent with a non-monotonic compression process.
When the model reaches a temporary PPL plateau---visible as a brief stabilization or decrease in uncurated PPL---the training signal weakens and the embedding space briefly decompresses.
This transient decompression restores partial alignment between geometric and content diversity, causing $\ratio$ to temporarily recover.
However, subsequent training quickly re-compresses the space, and the overall downward trend resumes.

\paragraph{Internal selection fails to prevent and may worsen collapse.}
RGDS produces a 13.0\% mean PPL excess over uncurated self-training ($p{=}0.02$, Cohen's $d \approx 1.8$) and a 30\% mean rare $n$-gram deficit ($p{=}0.04$).
While the effect sizes are statistically significant, we note that the rare $n$-gram $p$-value is marginal ($0.04$) and the absolute PPL excess, while consistent, is moderate.
These results support the claim that internal diversity selection \emph{fails to prevent and may worsen} collapse, but we stop short of claiming a dramatic acceleration effect.

\paragraph{The role of external embeddings.}
Frozen external embeddings (Contriever) maintain high selection ratios ($0.99 \pm 0.01$) because they are decoupled from the training feedback loop.
However, external selection improves PPL without improving rare $n$-gram survival, revealing that even correct geometric diversity selection does not address lexical tail loss.
This suggests that vocabulary collapse is driven by the generation process itself---the model stops \emph{producing} rare tokens---rather than by the selection process discarding them.

\paragraph{Practical implications.}
Our results have immediate practical relevance for iterative self-training pipelines:
(1) do not use the model's own embeddings for diversity-based data curation during iterative training;
(2) if embedding-based selection is desired, use frozen external embeddings, but do not expect this to preserve lexical diversity;
(3) do not rely on RankMe or $k$-NN probes to monitor diversity collapse---content-level metrics (rare $n$-grams, type-token ratio) are necessary.

\paragraph{Limitations.}
Our study is limited to GPT-2 scale models (125M and 355M parameters).
While the alignment gap mechanism---embedding space contraction under iterative training---is architecture-general, we have not verified it at billion-parameter scale.
The OpenWebText experiment uses a single seed, limiting cross-dataset statistical claims.
We use Contriever as the sole external embedding model; other encoders may exhibit different behavior.
Finally, our re-anchoring mitigation (Appendix~\ref{app:reanchor}) restores selection ratio but does not improve downstream metrics, leaving the development of effective countermeasures to future work.

\section{Conclusion}

We have identified the embedding alignment gap: a systematic divergence between internal embedding diversity and content diversity that emerges during iterative self-training.
This gap causes diversity-based data selection using internal representations to fail to prevent---and in some cases worsen---model collapse.
Our experiments reveal two measurement dissociation phenomena that challenge standard monitoring practices: perplexity and rare $n$-gram frequency decouple under diversity selection, and RankMe fails to detect ongoing diversity collapse.
These findings hold across model scales, datasets, and random seeds, suggesting the alignment gap is a fundamental property of iterative self-training rather than an artifact of specific configurations.
The field needs content-aware selection methods and diversity metrics that directly track lexical and distributional tails.

\appendix

\section{Cross-Seed Full Trajectories}
\label{app:seeds}

Table~\ref{tab:seed123} presents the RGDS selection ratio trajectory for seed=123.
The overall trend mirrors seed=42, with a terminal value of 0.78 and a rebound at Gen~5 (vs.\ Gen~4 for seed=42).
Seed=456 shows a terminal value of 0.80 with a rebound at Gen~6.
Full trajectories are provided below.

\begin{table}[h]
\centering
\small
\begin{tabular}{lcccccc}
\toprule
& \multicolumn{2}{c}{Seed 42} & \multicolumn{2}{c}{Seed 123} & \multicolumn{2}{c}{Seed 456} \\
\cmidrule(lr){2-3} \cmidrule(lr){4-5} \cmidrule(lr){6-7}
Gen & $\ratio$ & PPL & $\ratio$ & PPL & $\ratio$ & PPL \\
\midrule
0 & 1.03 & 37.2 & 1.01 & --- & 1.05 & --- \\
1 & 0.98 & 41.8 & 0.97 & --- & 1.00 & --- \\
2 & 0.93 & 46.5 & 0.94 & --- & 0.96 & --- \\
3 & 0.87 & 50.1 & 0.89 & --- & 0.90 & --- \\
4 & 0.91$^\dagger$ & 51.8 & 0.86 & --- & 0.87 & --- \\
5 & 0.84 & 53.2 & 0.88$^\dagger$ & --- & 0.83 & --- \\
6 & 0.80 & 53.8 & 0.82 & --- & 0.85$^\dagger$ & --- \\
7 & 0.76 & 54.2 & 0.78 & 55.1 & 0.80 & 53.8 \\
\bottomrule
\end{tabular}
\caption{RGDS selection ratio and PPL across 3 seeds (GPT-2 125M, WikiText-103). $^\dagger$ marks transient rebounds. Terminal Gen~7 metrics are available for all seeds; intermediate PPL is shown for seed=42 only (the primary experiment).}
\label{tab:seed123}
\end{table}

\section{Embedding Re-anchoring}
\label{app:reanchor}

As a simple mitigation, we test \emph{embedding re-anchoring}: every $K{=}3$ generations, the selection embeddings are recomputed using the original pretrained weights rather than the current model.
This resets the geometric diversity signal to its pre-collapse state.

Results (Table~\ref{tab:reanchor}) show that re-anchoring successfully restores selection ratio: the terminal value is 0.93 (vs.\ 0.76 for RGDS), with a sawtooth pattern as ratio drops within each $K$-step cycle then recovers at re-anchoring points.
However, this ratio recovery does not translate to meaningful downstream improvement: PPL improves by only 4.4\% over RGDS ($51.8$ vs $54.2$; $p{=}0.12$, not significant) and rare $n$-gram improves by 8.3\% ($0.0065$ vs $0.006$; $p{=}0.31$, not significant).

\begin{table}[h]
\centering
\small
\begin{tabular}{lccc}
\toprule
& RGDS & Re-anch. & Unc. \\
\midrule
$\ratio$ (Gen 7) & 0.76 & 0.93 & --- \\
PPL (Gen 7) & 54.2 & 51.8 & 48.0 \\
Rare (Gen 7) & .006 & .0065 & .009 \\
Rank corr. & $-$.090 & $-$.082 & $-$.085 \\
\midrule
$\ratio$ $p$-value & \multicolumn{2}{c}{\textbf{sig.}} & --- \\
PPL $p$-value & \multicolumn{2}{c}{0.12 (n.s.)} & --- \\
Rare $p$-value & \multicolumn{2}{c}{0.31 (n.s.)} & --- \\
\bottomrule
\end{tabular}
\caption{Re-anchoring vs.\ RGDS vs.\ Uncurated (GPT-2 125M, seed=42). Re-anchoring restores selection ratio but does not significantly improve downstream metrics, suggesting the root cause lies in the generation process rather than the selection process.}
\label{tab:reanchor}
\end{table}

This result is informative: it confirms that the alignment gap is a real and diagnosable phenomenon (the ratio responds to intervention), but reveals that correcting the selection signal alone is insufficient.
The generation process itself has already drifted toward homogeneity; selecting ``truly diverse'' data from a homogeneous pool cannot restore what the model has stopped producing.

\begin{thebibliography}{20}

\bibitem[Abbas et~al.(2023)]{abbas2023semdedup}
Abbas, M., Tirumala, K., Simig, D., Ganguli, S., and Morcos, A.~S.
\newblock {SemDeDup}: Data-efficient learning at web-scale through semantic deduplication.
\newblock In \emph{Proceedings of ACL}, 2023.

\bibitem[Alemohammad et~al.(2024)]{alemohammad2024selfconsuming}
Alemohammad, S., Casco-Rodriguez, J., Luzi, L., Humayun, A.~I., Babaei, H., LeJeune, D., Siahkoohi, A., and Baraniuk, R.~G.
\newblock Self-consuming generative models go {MAD}.
\newblock In \emph{Proceedings of ICLR}, 2024.

\bibitem[Dohmatob et~al.(2024)]{dohmatob2024tale}
Dohmatob, E., Feng, Y., and Gu, Q.
\newblock A tale of tails: Model collapse as a change of scaling laws.
\newblock In \emph{Proceedings of ICML}, 2024.

\bibitem[Garrido et~al.(2023)]{garrido2023rankme}
Garrido, Q., Balestriero, R., Najman, L., and LeCun, Y.
\newblock {RankMe}: Assessing the downstream performance of pretrained self-supervised representations by their rank.
\newblock In \emph{Proceedings of ICML}, 2023.

\bibitem[Izacard et~al.(2022)]{izacard2022contriever}
Izacard, G., Caron, M., Hosseini, L., Riedel, S., Bojanowski, P., Joulin, A., and Grave, E.
\newblock Unsupervised dense information retrieval with contrastive learning.
\newblock \emph{Transactions on Machine Learning Research}, 2022.

\bibitem[Merity et~al.(2016)]{merity2016pointer}
Merity, S., Xiong, C., Bradbury, J., and Socher, R.
\newblock Pointer sentinel mixture models.
\newblock \emph{arXiv preprint arXiv:1609.07843}, 2016.

\bibitem[Papyan et~al.(2020)]{papyan2020prevalence}
Papyan, V., Han, X.~Y., and Donoho, D.~L.
\newblock Prevalence of neural collapse during the terminal phase of deep learning training.
\newblock \emph{Proceedings of the National Academy of Sciences}, 117(40):24652--24663, 2020.

\bibitem[Shumailov et~al.(2024)]{shumailov2024model}
Shumailov, I., Shumaylov, Z., Zhao, Y., Gal, Y., Papernot, N., and Anderson, R.
\newblock The curse of recursion: Training on generated data makes models forget.
\newblock \emph{Nature}, 2024.

\bibitem[Tirumala et~al.(2023)]{tirumala2023d4}
Tirumala, K., Simig, D., Aghajanyan, A., and Morcos, A.~S.
\newblock {D4}: Improving {LLM} pretraining via document de-duplication and diversification.
\newblock \emph{arXiv preprint arXiv:2308.12284}, 2023.

\bibitem[Xie et~al.(2023)]{xie2023dsir}
Xie, S.~M., Santurkar, S., Ma, T., and Liang, P.
\newblock Data selection for language models via importance resampling.
\newblock In \emph{Proceedings of NeurIPS}, 2023.

\end{thebibliography}

\end{document}
